\subsection{S-Movies}

Skinsoft propose la gestion de collections cinématographiques, avec un espace de travail dédié, S-Movies. C'est au sein de cet espace que l'entreprise souhaite mettre en place l'indexation automatique des items films. 

Cet espace de travail repose sur la description des collections en différents niveaux du FRBR comme recommandé par la Fédéréation Internationale des Archives du Film. En effet, la FIAF recommande quatre niveaux de description hiérarchique pour le catalogage des films : l'oeuvre/expression, qui correspond à l'idée de l'oeuvre originale, la variante, qui est, comme son nom l'idique, une version différente de l'oeuvre originale, la manifestation, qui correspond à la diffusion de l'oeuvre, sa sortie, l'item, qui correspond à l'exemplaire physique ou numérique de l'oeuvre. Une oeuvre peut avoir plusieurs diffusion, et ainsi être liée à plusieurs manifestations. De plus, la collection peut compter plusieurs exemplaires physiques ou numériques d'une seule oeuvre. \footcite{zotero-item-646}

%Expliquer la fusion des deux niveaux : oeuvre et expression adapté au contexte cinématographique, et ajout du niveau variante. 

Ainsi pour une même oeuvre, il peut exister diverses diffusions/manifestations, et des collections aux formes multiples, au format pellicule, DVD, numérique... Une oeuvre originale peut elle-même subir des variantes : versions restaurées, montages différents, doublages et dialogues différents. Les niveaux du FRBR recomandés par la FIAF permettent alors de distinguer ces niveaux et les différences qui existent pour chacun de ces niveaux afin de mieux les décrire. Ce mode de description permet de séparer la description intelectuelle de la description matérielle. 

Toutefois, ce découpage en différents niveaux n'est pas toujours intuitif ni facilement déployable dans une application de gestion de collections. Les termes films et copies employés entre professionnels sont remplacés par des termes hiérarchiques opaques tels que manifestation, non spécifique à l'audiovisuel. (Oeuvre non plus n'est pas spécifique à l'audiovisuel)
Il n'est pas toujours évident de distinguer la variante d'une oeuvre. Où se trouve la limite entre l'oeuvre et sa variante ? La FIAF donne des pistes et des cas d'application mais ça ne couvre pas la complexité des collections réelle. 

%mentionner des exemples concrets (voir notes de stage) donner les exemples que donne la fiaf et montrer que ça ne s'applique pas à toutes les collections audiovisuelles


Une manifestation peut être par exemple une pré diffusion, une distribution en salles, une télédiffusion, une diffusion en ligne. \footcite{zotero-item-646}
Toutefois, les institutions qui conservent des archives audiovisuelles non diffusées, qui ne sont pas vraiment des oeuvres mais seulement des rushes documentaires par exemple, s'adaptent très mal au FRBR. La description se fait généralement au niveau de l'item. 

De plus, le modèle FRBR a été conçu pour les bilbliothèques au départ. Ainsi, le terme de "manifestation" par exemple parait opaque appliqué à des collections audiovisuelles. La FIAF tente d'adapter ce modèle aux collections audiovisuelles dans son manuel, mais les frontières entre les niveaux peuvent parfois être fines. Chaque institution choisit les niveaux qu'elle souhaite mettre en place au sein de ses pratiques de catalogage. Chez les clients de skinsoft, le niveau variante, par exemple reste peu utilisé puisque la frontière entre une oeuvre et sa variante sont complexes à définir. 

Au sein de l'application, dans l'espace S-Movies, la hiérarchisaton n'est pas intuitive. 

Lien avec l'indexation : crée des éléments qui se rattacheront à l'oeuvre et à l'item. 
Les métadonnées seront rattachées à l'item numérique (cf specification sur item/media).

Certains clients chez Skinsoft ont déjà une pratique manuelle de segmentation de leurs séquences audiovisuelles. Une institution en particulier qui préserve des archives films muettes documentaires, possède déjà des marqueurs de séparation de séquences avec les timecodes correspondant dans sa base de données. Dans ce cas, la segmentation automatique viendrait être entraînée sur ces corpus et pourrait remplacer cette pratique manuelle pour les archives non annotées. 

Problématique : les métadonnées issues de la segmentation automatique seraient liées à l'item numérique, segmenté en plusieurs médias qui correspondraient chacun à des séquences. Mais ces métadonnées décrivent le contenu sémantique de l'oeuvre. Les métadonnées doivent ainsi pouvoir être reversées sur la notice de l'oeuvre. 

Avant de pouvoir mettre en place la segmentation automatique, beaucoup de problèmes résident dans les entremêlements de liens entre les niveaux hiérarchiques du FRBR. Il faudrait pouvoir régler les confusions sur les liens entre les notices dans l'application avant de mettre en place une autre fonctionnalité. Sans avoir l'ambition de régler ces problèmes, nous proposerons ici toutefois une mise en place d'un processus de segmentation adapté et conforme au modèle FRBR. 

spécification sur l'item numérique : 
un item numérique est indexé dans la liste des objets de collection, et visible dans la liste des médias. les deux notices otn des champs partagés et des champs complémentaires. dans le cas d'une architecture fondée sur le FRBR, l'objet numérique est lié à une manifestation, une oeuvre. 
l'application doit permettre la création automatique d'un item numérique à l'ingest d'un fichier numérique, la lecture depuis l'item des données portées par le média/fichier, la lecture depuis le média des données portées par l'item. 
le fichier numérique formant un tout avec l'objet numérique est consultable dans le DAM avec possibilté de lire les données du FRBR
réutiliser sur le média le même métamasque que celui de l'objet. quand on paramètre le menu d'une notice média, card spécifique dédié pour apeller les métamasques de l'item. 
sur l'entité média je dois pouvoir utiliser les champs des cibles suivantes : média et notice de média, item, manifestation, oeuvre, 

à intégrer : 
lien entre le fichier numérique brut et la notice descriptive 
fichier numérique : métadonnées techniques
notice : métadonnées sémantiques 
ici fusion : le fichier numériue et la notice forment un tout 
facilite la mise en place de l'IA : analyse du fichier numérique et versement des métadonnées dans sa notice descriptive.